% 19 November 2026. Drafted for Edward Baker with OpenAI Codex assistance.
% AGC-E02; dispatched configuration gpt-6.1-sol, effort ultra.
% Dispatch metadata is recorded, not independently attested backend metadata.
\section{Integer pole powers and a single complex exponent}
\label{sec:hierarchy}

The natural indices of the rational hierarchy are nonnegative integers
attached separately to the two target poles. Their finite algebra is
proved first. A different construction later extends the exponent of
the \emph{combined} quadratic denominator to one complex variable.
The distinction is essential: no recurrence with two independently
complex split indices is established here
\cite{VCDCMManuscript,VCDCMDiscussion08}.

\subsection{Rational traces and a legitimate infinite-limit passage}

On a counted pair chart use $w_j=\rho_j-c$, $Q_j=w_j^2-D$ for
each exterior occurrence. Strict exterior separation on compact
subcharts, together with \eqref{agc:pair-count}, makes
\begin{equation}
 S_n=\sum_jQ_j^{-n}\quad(n\ge1),\qquad
 T_n=\sum_jw_jQ_j^{-n}\quad(n\ge2),\qquad M=S_1
 \label{agc:hier-ST}
\end{equation}
absolute. Their respective large-$|\rho|$ tails at fixed orders are
$O(N^{1-2n}\log N)$ and $O(N^{2-2n}\log N)$.
The first odd order $T_1$ is deliberately absent.

For either weight $k(z)=Q_*^{-n}$ or
$k(z)=(z-c)Q_*^{-n}$ in these ranges,
\begin{equation}
 \sum_j k(\rho_j)=-\frac1{2\pi\ii}
                       \oint_\Gamma k(z)B(t,z)\,\dd z.
 \label{agc:hier-trace}
\end{equation}
Indeed, for one exterior pole $k(z)/(z-\rho)$ the sum of residues
at the target poles is $-k(\rho)$, since there is no residue at
infinity. Product regularizer constants multiply $\oint k=0$.
The target contribution $kQ_{*,z}/Q_*$ has all its poles enclosed
and decays faster than $z^{-1}$, so its combined residue also vanishes.
This proves \eqref{agc:hier-trace} by locally uniform product
convergence. The contour must enclose both targets; separating their
residues can introduce nonzero terms which cancel only in the sum.

Here is the precise passage used for the hierarchy. At a reference
slice, truncate the opposite-paired product with all conjugates,
ties and multiplicities retained, and include the target pair. Give
each resulting polynomial its own heat evolution initialized at that
slice. Its spatial jets converge on the cleared contour. Its first
time logarithmic jet is determined there by
\begin{equation}
 B_t=-\frac{H_{zzz}}H+\frac{H_zH_{zz}}{H^2}
     =-B_{zz}-2BB_z.
 \label{agc:hier-Burgers}
\end{equation}
Differentiating the fixed-contour coordinate and trace formulas,
including the moving kernels, therefore passes their first time jets
to the limit. At a finite multiple-root slice one first perturbs the
factors to simple roots and uses the same contour argument. Thus a
finite polynomial identity passes to the complete field once every
nonabsolute first-power sum has cancelled and only fixed-order
absolute moment products remain. This slice argument asserts no
common time interval for the polynomial approximants and no
uniformity as their moment order grows.

\subsection{The integer split recurrence}

Where $D\ne0$ choose a local analytic $s=\sqrt D$, and put
$h=2s$, $X_j=w_j-s$, $Y_j=w_j+s=X_j+h$. Define
\begin{equation}
 P_{p,q}=\sum_jX_j^{-p}Y_j^{-q},\qquad
 p,q\in\mathbb Z_{\ge0},\quad p+q\ge2.
 \label{agc:hier-P}
\end{equation}
These are absolute even on the axes. Their tails are
$O(N^{1-p-q}\log N)$ at fixed indices and fixed cleared targets.
For $m,n\ge1$ write
\[
 C_{m,n}=\sum_{i=2}^{m}P_{i,0}P_{m+1-i,n}
          +\sum_{j=1}^{n-1}P_{m,j}P_{0,n+1-j},
\]
and on an axis use
$C_{m,0}=\sum_{i=2}^{m-1}P_{i,0}P_{m+1-i,0}$ and
$C_{0,n}=\sum_{j=2}^{n-1}P_{0,j}P_{0,n+1-j}$.
An empty sum is zero. With $A_h=2/h-hM$, the recurrence is
\begin{align}
 P_{p,q}'={}&p\{C_{p+1,q}-(p+3)P_{p+2,q}
                  -(q+2)P_{p+1,q+1}+A_hP_{p+1,q}\}\notag\\
 &+q\{C_{p,q+1}-(p+2)P_{p+1,q+1}
                  -(q+3)P_{p,q+2}-A_hP_{p,q+1}\},
          &&p,q\ge1, \label{agc:hier-mixed}\\
 P_{p,0}'={}&p\{C_{p+1,0}-(p+3)P_{p+2,0}-2P_{p+1,1}
                                     +(2/h)P_{p+1,0}\},
          &&p\ge2, \label{agc:hier-axis-one}\\
 P_{0,q}'={}&q\{C_{0,q+1}-2P_{1,q+1}-(q+3)P_{0,q+2}
                                     -(2/h)P_{0,q+1}\},
          &&q\ge2. \label{agc:hier-axis-two}
\end{align}
Also $h'=4/h-2hM$. These identities hold through exterior
multiplicities on separated target charts.

\begin{proof}
Work first with finitely many simple roots. The target velocities are
$z_+'=2/h-2P_{1,0}$, $z_-'=-2/h-2P_{0,1}$, where the first-order
symbols are used only at this finite stage. Since
$X_j^{-1}-Y_j^{-1}=h/Q_j$, their difference gives $h'$.
An exterior root receives the target force
$2(X_j^{-1}+Y_j^{-1})$. For
$f_{m,n}(w)=(w-s)^{-m}(w+s)^{-n}$, the exterior--exterior force
contributes
\[
 U_{m,n}=-2\sum_{j<k}
             \frac{f_{m,n}(w_j)-f_{m,n}(w_k)}{w_j-w_k}
\]
to the block multiplied by $p$ or $q$ after differentiating its
weight. Reciprocal-power subtraction gives the finite identity
\begin{equation}
 -\frac{f_{m,n}(u)-f_{m,n}(v)}{u-v}
 =\sum_{i=1}^{m}f_{i,0}(u)f_{m+1-i,n}(v)
   +\sum_{j=1}^{n}f_{m,j}(u)f_{0,n+1-j}(v).
 \label{agc:hier-split-DD}
\end{equation}
To verify it, subtract the first factors while holding
$(v+s)^{-n}$ fixed, then subtract the second factors while
holding $(u-s)^{-m}$ fixed. The full left side is symmetric in
$u,v$, although this decomposition need not look symmetric.
Thus $U_{m,n}$ is the ordered off-diagonal sum on the right.
Its diagonal subtractions are $-mP_{m+1,n}-nP_{m,n+1}$.
For $m,n\ge1$, its two endpoint products are
$(P_{1,0}+P_{0,1})P_{m,n}$, and its remaining products are $C_{m,n}$.

The differentiated first block is
\[
 p\{(2/h-2P_{1,0})P_{p+1,q}
                  -2P_{p+2,q}-2P_{p+1,q+1}+U_{p+1,q}\};
\]
the second is the same expression with prefactor $q$, target
velocity $-2/h-2P_{0,1}$, and raised second index. Adding the
off-diagonal formulas cancels the common first-order sums; their
remaining difference is $P_{1,0}-P_{0,1}=hM$. The diagonals and
the two explicit target-force terms give exactly
\eqref{agc:hier-mixed}. On an axis both endpoint products contain
the same first-order sum and cancel it completely, leaving
\eqref{agc:hier-axis-one}--\eqref{agc:hier-axis-two}.
After these cancellations every surviving series is absolute.
The weight in \eqref{agc:hier-P} has the same fixed-contour trace
as \eqref{agc:hier-trace}. The preceding slice-jet passage therefore
proves the complete equations, including exterior multiplicities.
\end{proof}

Rows attached to different adjacent edges are constrained by their
shared divisor. If $P_{p,q}^{(n)}$ denotes \eqref{agc:hier-P} for
$(x_n,x_{n+1})$, restoring the omitted shared neighbor gives, for $r\ge2$,
\begin{equation}
 P_{r,0}^{(n)}+(-g_n)^{-r}
       =P_{0,r}^{(n+1)}+g_{n+1}^{-r}.
 \label{agc:hier-overlap}
\end{equation}
Both sides equal $\sum_{k\ne n+1}(x_k-x_{n+1})^{-r}$.
Within one row, $Y-X=h$ also gives
$hP_{p,q}=P_{p,q-1}-P_{p-1,q}$ for $p,q\ge1$,
$p+q\ge3$. At $(1,1)$ the difference must be combined before
summation. These necessary compatibility identities do not assert
that an arbitrary moment array comes from one divisor.

\subsection{Symmetric variables and their collision-regular evolution}

The split labels exchange when $s$ is replaced by $-s$. Their
regular symmetric form is \eqref{agc:hier-ST}, since
\begin{equation}
 S_n=P_{n,n},\qquad
 T_n=\tfrac12(P_{n-1,n}+P_{n,n-1}).
 \label{agc:hier-change}
\end{equation}
Conversely, put $r=\max(p,q)$ and write
\[
 X^{-p}Y^{-q}
     =\frac{(w-s)^{r-p}(w+s)^{r-q}}{(w^2-D)^r}.
\]
Its numerator has degree at most $2r-2$. Replacing even powers
$w^{2k}$ by $(Q+D)^k$ and odd powers by $w(Q+D)^k$ expresses
the result using only $S_n$ for $n\ge1$ and $T_n$ for $n\ge2$,
with polynomial coefficients in $s$. This proves equivalence of the
integer descriptions in their convergent ranges.

The collision-regular equations are
\begin{align}
 S_n'={}&2n\sum_{k=1}^{n}S_kS_{n+1-k}
       +2n\sum_{k=2}^{n}(T_kT_{n+2-k}+DS_kS_{n+2-k})\notag\\
       &-4n(n+2)S_{n+1}-4n(n+3)DS_{n+2},&&n\ge1,
                                                  \label{agc:hier-S}\\
 T_n'={}&2(2n-1)\sum_{k=2}^{n}T_kS_{n+1-k}
       +4nD\sum_{k=2}^{n}T_kS_{n+2-k}\notag\\
       &-4(n^2+n-1)T_{n+1}-4n(n+3)DT_{n+2},&&n\ge2.
                                                  \label{agc:hier-T}
\end{align}

For completeness, the symmetrization giving these coefficients can
be checked without manipulating a singular split label. At a finite
simple slice put $P=T_1=\sum_jw_j/Q_j$. The relative root equation is
\begin{equation}
 w_j'=4w_j/Q_j+2P+2\sum_{k\ne j}(w_j-w_k)^{-1}.
 \label{agc:hier-relative}
\end{equation}
For $A=(u^2-D)^{-1}$, $B_0=(v^2-D)^{-1}$, and
$L_r=\sum_{i=0}^{r-1}A^{r-1-i}B_0^i$ with $L_0=0$, elementary
subtraction gives
\begin{equation}
 \begin{split}
 \frac{A^r-B_0^r}{u-v}&=-(u+v)AB_0L_r,\\
 \frac{uA^r-vB_0^r}{u-v}
       &=-AB_0L_{r-1}-(D+uv)AB_0L_r.
 \end{split}
 \label{agc:hier-regular-DD}
\end{equation}
Here $AB_0L_r=\sum_{k=1}^r A^kB_0^{r+1-k}$. Summing over
unordered exterior pairs and subtracting the diagonals therefore gives
\begin{align*}
 -\sum_{i<j}\frac{w_iQ_i^{-n-1}-w_jQ_j^{-n-1}}{w_i-w_j}
 ={}&\frac12\sum_{k=1}^{n}(S_kS_{n+1-k}-S_{n+1})\\
 &+\frac D2\sum_{k=1}^{n+1}(S_kS_{n+2-k}-S_{n+2})\\
 &+\frac12\sum_{k=1}^{n+1}
                (T_kT_{n+2-k}-S_{n+1}-DS_{n+2}),\\
 -\sum_{i<j}\frac{Q_i^{-r}-Q_j^{-r}}{w_i-w_j}
 ={}&\sum_{k=1}^r(T_kS_{r+1-k}-T_{r+1}).
\end{align*}
In these two finite identities only, $T_1=P$ is allowed.
Differentiating $Q^{-n}$ using \eqref{agc:hier-relative} gives
\[
 nD'S_{n+1}-8n(S_{n+1}+DS_{n+2})-4nPT_{n+1}
       -4n\sum_{i<j}\frac{w_iQ_i^{-n-1}-w_jQ_j^{-n-1}}{w_i-w_j}.
\]
The $k=1,n+1$ odd endpoint products cancel $P$, and the
$D$-weighted endpoints cancel $-4nDMS_{n+1}$ from $nD'S_{n+1}$.
The diagonal coefficients left are $-4n(n+2)$ and $-4n(n+3)$,
proving \eqref{agc:hier-S}. For the odd weight use
\[
 (wQ^{-n})_w=(1-2n)Q^{-n}-2nDQ^{-n-1},\qquad
 (wQ^{-n})_D=nwQ^{-n-1}.
\]
The second divided-difference sum then cancels both terms involving
$P$; its $D$-weighted last endpoint cancels the interaction in
$nD'T_{n+1}$. The remaining direct and diagonal coefficients are
$-4(n^2+n-1)$ and $-4n(n+3)$, giving \eqref{agc:hier-T}.
The fixed-order slice passage proves both complete laws, through
target and exterior multiplicities.

In particular,
\begin{equation}
 \begin{split}
 M'&=2M^2-12S_2-16DS_3,\\
 S_2'&=8MS_2-32S_3+4DS_2^2-40DS_4+4T_2^2,\\
 T_2'&=(6M+8DS_2)T_2-20T_3-40DT_4.
 \end{split}
 \label{agc:hier-low}
\end{equation}
The spatial term $4T_2^2$ has a positive sign. All signs refer to
forward physical time; backward transport must reverse the time
derivative of the whole equation. These are componentwise identities
of the given complete field. They do not prove well-posedness for
arbitrary moment arrays or stability after a finite-order closure.

At a target merger the regular moments remain analytic, but a raw
split coordinate generally has the form $A(t)+s(t)C(t)$, with
$A,C$ analytic. Its coefficient at the double is
$C(t_*)=(p-q)\sum_jw_j(t_*)^{-p-q-1}$.
Since $D'(t_*)=2$, a nonzero coefficient produces a derivative of
order $1/s$. Thus collision regularity belongs to the symmetric
description, not to each separately labelled pole coordinate.

\subsection{The full inverse-power family and its Mellin meaning}

Choose a fixed length $\ell>0$. For each exterior occurrence select
a spatial logarithm of $Q(t,w)=w^2-D(t)$ near its location and set
\begin{equation}
 f_p(t,w)=\exp\{p(2\log\ell-\operatorname{Log}Q(t,w))\},\qquad
 Z(p,t)=\sum_j f_p(t,w_j),\qquad \Re p>\tfrac12.
 \label{agc:hier-Z}
\end{equation}
Then $Z(n,t)=\ell^{2n}S_n(t)$ for integer $n\ge1$.
The infinite inverse-power family retains all denominator scales;
increasing a real order emphasizes smaller exterior denominators.
For positive $Q_1<Q_2$, the ratio of their order-$n$ contributions
is $(Q_2/Q_1)^n$. This compares contributions within one order,
not dimensionful values of different moments. Nonreal denominators
carry phases and need not provide a positive moment sequence.

When every $Q_j>0$, define the denominator heat trace
$\Theta(v,t)=\sum_j\exp(-vQ_j/\ell^2)$ for $v>0$. The gamma
integral gives the genuine Mellin representation
\begin{equation}
 Z(p,t)=\frac1{\Gamma(p)}\int_0^\infty
                         v^{p-1}\Theta(v,t)\,\dd v,
                         \qquad \Re p>\tfrac12.
 \label{agc:hier-Mellin}
\end{equation}
Indeed, for $a=\Re p>1/2$, the integral of the sum of absolute
values is $\Gamma(a)\sum_j(\ell^2/Q_j)^a<\infty$.
Tonelli followed by Fubini proves the formula. It also holds with
the compatible logarithm if all $Q_j$ lie in a common closed sector
$\Re Q_j\ge\eta|Q_j|>0$, since the same argument uses
$(\Re Q_j)^{-a}$. A general complex exterior need not satisfy this
sector condition. In that case \eqref{agc:hier-Z} is the primary
complex-power trace; an unrotated positive-ray Laplace integral
must not be asserted for it.

\subsection{Local logarithms and the complete complex-time law}

For this analytic extension work at real times with a
conjugation-invariant pair, so $c,D$ are real. Assume, in addition to the complete
product, a fixed strip $|\Im\rho|\le A_0$, the projected unit count
$\#\{\rho:v\le\Re\rho\le v+1\}\le C\log(2+|v|)$,
and safe-line bounds
$|B(t,x\pm\ii Y)|\le C\log(2+|x|)$ with $Y>A_0$.
These are explicit source inputs, not consequences of a moment
recurrence. Keep $c,D$ in a compact counted chart and choose a large
centered cutoff $R_0$ so that $\Re(w^2-D)\ge(\Re w)^2/2$ for
$|\Re w|\ge R_0$, $|\Im w|\le Y$. Use the principal logarithm
on those two half-strips. Enclose the finitely many central exterior
clusters by disjoint cleared disks avoiding $Q=0$. On each disk
choose one analytic logarithm of $Q$ and continue it in time;
shrinking the chart makes this possible even through an exterior
merger. There is a chart constant $\Theta_0$ bounding the absolute
imaginary parts of these logarithms.

Let $\mathcal S$ denote $(2\pi\ii)^{-1}$ times integration over
the positive central disk boundaries and the boundaries of the two
infinite half-strips. On the latter the lower rays run rightward,
the upper rays leftward, the finite right connector downward and
the finite left connector upward. The central target pair is excluded.
For $h(t,w)=B(t,c(t)+w)$, the residue theorem gives
\begin{equation}
 \begin{split}
 Z&=\mathcal S(f_ph),\\
 Z_t&=\mathcal S((f_p)_t h+f_ph_t)\\
    &=\mathcal S\big(\{(f_p)_t-(f_p)_{ww}-c'(f_p)_w\}h
                                      +(f_p)_wh^2\big),\\
 h_t&=-h_{ww}-2hh_w+c'h_w,\qquad
 (f_p)_t=pD'f_p/Q.
 \end{split}
 \label{agc:hier-complex-time}
\end{equation}
The derivatives of $f_p,h$ here hold $w$ fixed.

We give the convergence details because a formal sum of zero
velocities at infinitely many multiple roots would not justify this
law. At fixed $1/2<b\le\Re p\le r$, counting and the chosen phases give
\begin{equation}
 |Z(p,t)|\le C_{b,r}\exp(\Theta_0|\Im p|),\qquad
 |Z_t(p,t)|\le C_{b,r}(1+|p|)^2\exp(\Theta_0|\Im p|)
 \label{agc:hier-p-envelope}
\end{equation}
on compact subcharts. For the value estimate, the tails are bounded
by $\sum_{n\ge1}n^{-2b}\log(2+n)$; central terms are finite and
separated. On the safe rays the absolute reciprocal-square and
reciprocal-cube sums yield $|h_w|+|h_{ww}|=O(\log(2+|x|))$.
Thus $h_t=O(\log^2(2+|x|))$ by \eqref{agc:hier-complex-time}.
Moreover
\[
 (f_p)_w=-2pw f_p/Q,\qquad
 (f_p)_{ww}=\{-2p/Q+4p(p+1)w^2/Q^2\}f_p.
\]
Both differentiated integrands consequently have an integrable
majorant $C(1+|p|)^2|x|^{-2b}\log^2(2+|x|)
\exp(\Theta_0|\Im p|)$. Cleared finite pieces have bounded
spatial jets and a positive $|Q|$ floor, giving the same exponent
envelope with a finite constant.

To remove the distant vertical closures, choose their real projections
in successive unit intervals at distance at least
$c/\log(2+|x|)$ from all zero projections. Such cuts exist by
deleting the short forbidden intervals and using the unit count.
Comparing $h$ on each cut with its value on the upper safe line
via the once-subtracted logarithmic derivative gives
$h=O(\log^2(2+|x|))$ there: the near unit shell pays inverse
projection clearance and the far shells have inverse-square decay.
The bounded-length connector integral is therefore
$O(|x|^{-2b}\log^2(2+|x|))$ and tends to zero. This proves the
spatial residue formula. The displayed majorants justify time
differentiation and normal convergence in $p$ for $\Re p>1/2$.
Finally, the difference between the two time integrands in
\eqref{agc:hier-complex-time} is the spatial derivative of
$-f_ph_w+(f_p)_wh-f_ph^2+c'f_ph$. Its circle integrals vanish,
its connector endpoint values cancel, and its values at infinite
ray endpoints tend to zero. Integration by parts proves the last
representation, including the full interaction term $(f_p)_wh^2$.

Two analytic time charts can differ on finitely many central clusters
by logarithm shifts $2\pi\ii k_C$. On their overlap,
\begin{equation}
 Z_B-Z_A=\sum_C(e^{-2\pi\ii k_Cp}-1)Z_{A,C}.
 \label{agc:hier-branch-transition}
\end{equation}
The same relation differentiates in time on the local cluster
contours. All corrections and their time derivatives vanish at
integer $p$. When real exterior denominators are positive, as for
a count-two disk containing the target segment, conjugate clusters
can be assigned conjugate logs, including opposite arguments
$+\pi,-\pi$ at nonreal conjugate locations with negative-real $Q$. This gives
$Z(p)=\overline{Z(\bar p)}$. If instead every negative-real
$Q_j$ is assigned $+\pi$ on a single slice, the exact defect is
\[
 Z_+(p)-\overline{Z_+(\bar p)}
   =-2\ii\ell^{2p}\sin(\pi p)\sum_{Q_j<0}|Q_j|^{-p}.
\]
The exceptional sum is finite. Resetting arguments at every time
need not define an analytic time chart. These branch restrictions
are essential away from the integer hierarchy.

\subsection{Damped recovery of the forcing}

The complex representation recovers the inherited forcing exactly.
Set $\delta=D/\ell^2$, $m=\ell^2M$, $\sigma=(t-t_0)/\ell^2$,
fix $\epsilon>0$, and use
\[
 K_\epsilon(p,t)=\frac{12e^{\epsilon(p-2)^2}}{p-2}
              +\frac{16\delta(t)e^{\epsilon(p-3)^2}}{p-3},
 \qquad G_\epsilon(p)=\frac{e^{\epsilon(p-1)^2}}{p-1}.
\]
For upward vertical lines $V_q$, define the two-line difference
$\mathcal V_{b,r}(F)=(2\pi\ii)^{-1}
(\int_{V_r}F\,\dd p-\int_{V_b}F\,\dd p)$. Then
\begin{equation}
 \begin{split}
 F(t)&=\mathcal V_{1,4}(K_\epsilon Z)
       =12Z(2,t)+16\delta Z(3,t)
       =\ell^4(12S_2+16DS_3),\\
 m(t)&=\mathcal V_{3/4,3/2}(G_\epsilon Z)=Z(1,t),\\
 \frac{\dd m}{\dd\sigma}&=2m^2-F,\qquad
 \frac{\dd\delta}{\dd\sigma}=2-4m\delta.
 \end{split}
 \label{agc:hier-forcing-contour}
\end{equation}
To prove this, close a counterclockwise rectangle between the chosen
verticals. The damping equals one at each residue and has modulus
$e^{\epsilon((\Re p-n)^2-(\Im p)^2)}$.
The horizontal integrals therefore vanish as their height tends
to infinity by \eqref{agc:hier-p-envelope}; the vertical integrals
are absolutely convergent. The same Gaussian dominates the time
derivative, including
$16(D'/\ell^2)e^{\epsilon(p-3)^2}Z/(p-3)$ from differentiating
the kernel. Residues give the first two formulas, and
\eqref{agc:pair-law}, \eqref{agc:hier-low} give the last two.
On a branch overlap the correction \eqref{agc:hier-branch-transition}
vanishes at the residue indices. The apparent poles of its product
with either kernel are removable, so the complete two-line values
and their derivatives agree across charts.

These contours reorganize the complete forcing. They provide neither
its sign nor an estimate stronger than evaluating the same integer
moments from the same information. In particular one cannot replace
the integer upper limits in \eqref{agc:hier-mixed} or
\eqref{agc:hier-S} by arbitrary complex indices. The single-complex-$p$
law \eqref{agc:hier-complex-time} is the proved analytic extension;
the independent canonical input needed to use the forcing for
continuation is considered next.
